

\bta{2004}

\begin{enumerate}
	%\renewcommand{\labelenumi}{\arabic{enumi}.}
	% A(\Alph) a(\alph) I(\Roman) i(\roman) 1(\arabic)
	%设定全局标号series=example	%引用全局变量resume=example
	%[topsep=-0.3em,parsep=-0.3em,itemsep=-0.3em,partopsep=-0.3em]
	%可使用leftmargin调整列表环境左边的空白长度 [leftmargin=0em]
	\item
（$30'$）粒子在一维无限深方势阱$V(x)$中运动，$V(x)=\begin{cases}\infty,&|x|>a\\0,&|x|<a\end{cases}$处于状态$\psi=\phi_1+\phi_3+2\phi_4$。这里$\phi_n,n=1,2,3,\cdots$是系统归一化的能量本征态。请问：
\begin{enumerate}
	%\renewcommand{\labelenumi}{\arabic{enumi}.}
	% A(\Alph) a(\alph) I(\Roman) i(\roman) 1(\arabic)
	%设定全局标号series=example	%引用全局变量resume=example
	%[topsep=-0.3em,parsep=-0.3em,itemsep=-0.3em,partopsep=-0.3em]
	%可使用leftmargin调整列表环境左边的空白长度 [leftmargin=0em]
	\item
粒子具有基态能量$E_1$几率；
\item 
粒子的平均能量（用基态能量$E_1$的倍数表示）；
\item 
态$\phi_4$中的节点数（在节点处，找到粒子的几率密度为零）；
\item 
态$\phi_3$的宇称。

\end{enumerate}

\banswer{
	原状态未归一化，其模方为 $1+1+4=6$，故归一化状态为
	\[
		|\psi\rangle=\frac{1}{\sqrt6}
		\left(|\phi_1\rangle+|\phi_3\rangle+2|\phi_4\rangle\right).
	\]
	\begin{enumerate}
		\item 测得基态能量 $E_1$ 的概率为 $P_1=1/6$。
		\item 对宽度为 $2a$ 的无限深势阱，$E_n=n^2E_1$，所以
		\[
			\langle H\rangle=\frac{E_1+E_3+4E_4}{6}
			=\frac{1+9+4\times16}{6}E_1=\frac{37}{3}E_1.
		\]
		\item $\phi_n$ 在势阱内部有 $n-1$ 个节点，故 $\phi_4$ 有 $3$ 个内部节点（不计两端边界）。
		\item 对称势阱的本征态宇称交替，$\phi_1$ 为偶态，故 $\phi_3$ 也为偶态。
	\end{enumerate}
}


\item 
考虑一维体系$\hat{H}=\frac{p^2}{2\mu}+V(x),\quad V(x)=V_0x^\lambda,\quad V_0>0,\quad \lambda=2,4,6\cdots$。设$\hat{H}$的本征波函数为$\psi_n$。
\begin{enumerate}
	%\renewcommand{\labelenumi}{\arabic{enumi}.}
	% A(\Alph) a(\alph) I(\Roman) i(\roman) 1(\arabic)
	%设定全局标号series=example	%引用全局变量resume=example
	%[topsep=-0.3em,parsep=-0.3em,itemsep=-0.3em,partopsep=-0.3em]
	%可使用leftmargin调整列表环境左边的空白长度 [leftmargin=0em]
	\item
证明动量在态$\psi_n$中的平均值为零；

\item 
求在态$\psi_n$中的动能平均值和势能平均值之间的关系。

	
\end{enumerate}

\banswer{
	\begin{enumerate}
		\item 束缚定态只差一个随时间变化的整体相位，故 $\langle x\rangle$ 不随时间变化。由 Ehrenfest 定理
		\[
			\frac{d\langle x\rangle}{dt}=\frac{\langle p\rangle}{\mu}=0,
		\]
		所以 $\langle p\rangle=0$。等价地，因势能为偶函数，本征态可取确定宇称，$\psi_n^*(-i\hbar\partial_x)\psi_n$ 为奇函数，其全空间积分为零。

		\item 对 $V(x)=V_0x^\lambda$，有 $xV'(x)=\lambda V(x)$。维里定理给出
		\[
			2\langle T\rangle=\langle xV'(x)\rangle
			=\lambda\langle V\rangle,
		\]
		故
		\[
			\boxed{\langle T\rangle=\frac{\lambda}{2}\langle V\rangle},
			\qquad
			\langle T\rangle=\frac{\lambda}{\lambda+2}E_n,
			\quad \langle V\rangle=\frac{2}{\lambda+2}E_n.
		\]
	\end{enumerate}
}



\item 
设归一化的状态波函数$|\psi\rangle$满足薛定鄂方程$i\hbar\frac{\partial}{\partial t}|\psi\rangle=\hat{H}|\psi\rangle$，定义密度算符（矩阵）为$\rho=|\psi\rangle\langle\psi|$。
\begin{enumerate}
	%\renewcommand{\labelenumi}{\arabic{enumi}.}
	% A(\Alph) a(\alph) I(\Roman) i(\roman) 1(\arabic)
	%设定全局标号series=example	%引用全局变量resume=example
	%[topsep=-0.3em,parsep=-0.3em,itemsep=-0.3em,partopsep=-0.3em]
	%可使用leftmargin调整列表环境左边的空白长度 [leftmargin=0em]
	\item
证明任意力学量$\hat{F}$在态$|\psi\rangle$中的平均值可表示为$Tr(\rho F)$；

\item 
求出$\rho$的本征值；

\item 
导出随时间演化的方程。
\end{enumerate}

\banswer{
	\begin{enumerate}
		\item 在任意完备正交基 $\{|i\rangle\}$ 中，
		\begin{align*}
			\operatorname{Tr}(\rho F)
			&=\sum_i\langle i|\psi\rangle\langle\psi|F|i\rangle
			=\langle\psi|F\left(\sum_i|i\rangle\langle i|\right)|\psi\rangle\\
			&=\langle\psi|F|\psi\rangle=\langle F\rangle.
		\end{align*}

		\item 由归一化条件 $\langle\psi|\psi\rangle=1$，
		\[
			\rho|\psi\rangle=|\psi\rangle,
		\]
		故沿 $|\psi\rangle$ 的本征值为 $1$。任意与 $|\psi\rangle$ 正交的态 $|u_\perp\rangle$ 满足 $\rho|u_\perp\rangle=0$，故其余本征值全为 $0$。

		\item 由薛定谔方程及其伴随式，
		\begin{align*}
			i\hbar\frac{\partial\rho}{\partial t}
			&=i\hbar\left(|\dot\psi\rangle\langle\psi|+|\psi\rangle\langle\dot\psi|\right)\\
			&=H|\psi\rangle\langle\psi|-|\psi\rangle\langle\psi|H
			=[H,\rho].
		\end{align*}
		这就是纯态密度算符的 von Neumann 方程。
	\end{enumerate}
}


\newpage
\item 
质量为$\mu$的粒子在三维各向同性谐振子势为$V(r)=\frac{kr^2}{2}=\frac{k(x^2+y^2+z^2)}{2}$中运动。求：
\begin{enumerate}
	%\renewcommand{\labelenumi}{\arabic{enumi}.}
	% A(\Alph) a(\alph) I(\Roman) i(\roman) 1(\arabic)
	%设定全局标号series=example	%引用全局变量resume=example
	%[topsep=-0.3em,parsep=-0.3em,itemsep=-0.3em,partopsep=-0.3em]
	%可使用leftmargin调整列表环境左边的空白长度 [leftmargin=0em]
	\item
第二激发态的能量；
\item 
第一激发态的简并度；
\item 
在基态中的不确定量$\Delta r\cdot\Delta p$，这里$\Delta r$是位置矢量的均方差根。定
义类似。
\end{enumerate}


\banswer{
	令 $\omega=\sqrt{k/\mu}$。三维各向同性谐振子的能量为
	\[
		E_N=\hbar\omega\left(N+\frac32\right),
		\qquad N=n_x+n_y+n_z=0,1,2,\ldots,
	\]
	其空间简并度为 $g_N=(N+1)(N+2)/2$。
	\begin{enumerate}
		\item 第二激发态对应 $N=2$，故 $E_2=\frac72\hbar\omega$。
		\item 第一激发态对应 $N=1$，故简并度 $g_1=3$。
		\item 基态中 $\langle\vec r\rangle=\langle\vec p\rangle=0$，且三个方向相互独立：
		\[
			\langle r^2\rangle=3\frac{\hbar}{2\mu\omega},
			\qquad
			\langle p^2\rangle=3\frac{\mu\hbar\omega}{2}.
		\]
		因此
		\[
			\Delta r=\sqrt{\frac{3\hbar}{2\mu\omega}},
			\qquad
			\Delta p=\sqrt{\frac{3\mu\hbar\omega}{2}},
			\qquad
			\boxed{\Delta r\,\Delta p=\frac{3\hbar}{2}}.
		\]
	\end{enumerate}
}


\item 
两个自旋均为 $1/2$ 的粒子 $1$ 和粒子 $2$ 组成的系统处于
$|\psi\rangle=a|0\rangle_1|1\rangle_2+b|1\rangle_1|0\rangle_2$，
其中 $|0\rangle$ 表示自旋朝下（沿 $-z$ 方向），$|1\rangle$ 表示自旋朝上。
当 $a$ 和 $b$ 均不为零时，此态不能写成两个单粒子态的直积，称为纠缠态。求上述纠缠态中：
\begin{enumerate}
	%\renewcommand{\labelenumi}{\arabic{enumi}.}
	% A(\Alph) a(\alph) I(\Roman) i(\roman) 1(\arabic)
	%设定全局标号series=example	%引用全局变量resume=example
	%[topsep=-0.3em,parsep=-0.3em,itemsep=-0.3em,partopsep=-0.3em]
	%可使用leftmargin调整列表环境左边的空白长度 [leftmargin=0em]
	\item
两个粒子的自旋互相平行的几率;
\item
两个粒子的自旋互相反平行的几率;
\item
此系统处于总自旋为 $\mathbf{0}$ 的几率;	
	
\item 
测量得到粒子 $ 1 $ 自旋朝上的几率多大; 发现粒子 $ 1 $ 自旋朝上时，粒子 $ 2 $ 处于什么 状态?


	
	
\end{enumerate}


\banswer{
	设状态已归一化，即 $|a|^2+|b|^2=1$。记 $|1\rangle=|\uparrow\rangle$、$|0\rangle=|\downarrow\rangle$。
	\begin{enumerate}
		\item 两个展开项中两粒子沿 $z$ 方向均反平行，故平行概率为 $0$。
		\item 反平行概率为 $|a|^2+|b|^2=1$。
		\item 总自旋单态为
		\[
			|0,0\rangle=\frac{|\uparrow\downarrow\rangle-|\downarrow\uparrow\rangle}{\sqrt2}.
		\]
		于是
		\[
			P_{S=0}=|\langle0,0|\psi\rangle|^2
			=\frac{|b-a|^2}{2}.
		\]
		\item 测得粒子 $1$ 自旋朝上的概率为 $|b|^2$；在得到该结果后，归一化的塌缩态为 $|\uparrow\rangle_1|\downarrow\rangle_2$，所以粒子 $2$ 必处于自旋朝下态 $|0\rangle_2$。
	\end{enumerate}
}

\item 
考虑到自旋轨道耦合的氢原子，其哈密顿量为 $\hat{H}=-\frac{\hbar^{2}}{2 \mu} \nabla^{2}+c(r)+\xi(r) \hat{\vec{L}} \cdot \hat{\vec{S}}$。
\begin{enumerate}
	%\renewcommand{\labelenumi}{\arabic{enumi}.}
	% A(\Alph) a(\alph) I(\Roman) i(\roman) 1(\arabic)
	%设定全局标号series=example	%引用全局变量resume=example
	%[topsep=-0.3em,parsep=-0.3em,itemsep=-0.3em,partopsep=-0.3em]
	%可使用leftmargin调整列表环境左边的空白长度 [leftmargin=0em]
	\item
	证明轨道角动量 $\hat{\vec{L}}$ 和 $\hat{\vec{S}}$ 不是此系统的守恒量, 而总角动量 $\hat{\vec{J}}=\hat{\vec{L}}+\hat{\vec{S}}$ 是守恒量。
	\item
	若自旋一轨道相互作用 $\xi(r) \hat{\vec{L}} \cdot \hat{\vec{S}}$ 可当作微扰，计算此系统基态能量的一级修正。 \\
	（$\hat{H}_{0}=-\frac{\hbar^{2}}{2 \mu} \nabla^{2}+V(r)$ 的本征能量为 $E_{n}^{(0)}$, 本征函数: $\psi_{n l m s}=R_{n l}(r) Y_{l m}(\theta, \varphi) \chi_{s}, \chi_{s}$为自旋波函数）
	
	
	
\end{enumerate}

\banswer{
	把哈密顿量写成 $H=H_0+\xi(r)\vec L\cdot\vec S$，其中球对称的 $H_0$ 与 $\vec L$、$\vec S$ 均对易。自旋--轨道项给出
	\[
		[\vec L\cdot\vec S,L_i]\neq0,
		\qquad
		[\vec L\cdot\vec S,S_i]\neq0,
	\]
	故 $\vec L$ 和 $\vec S$ 分别不守恒。另一方面
	\begin{align*}
		[\vec L\cdot\vec S,J_i]
		&=[\vec L\cdot\vec S,L_i]+[\vec L\cdot\vec S,S_i]=0,
	\end{align*}
	且 $[H_0,J_i]=0$，所以 $[H,\vec J]=0$，总角动量 $\vec J=\vec L+\vec S$ 守恒。

	氢原子基态为 $n=1,l=0$，因而
	\[
		\vec L\cdot\vec S
		=\frac12(J^2-L^2-S^2)=0.
	\]
	所以自旋--轨道微扰对基态的一级能量修正为
	\[
		\boxed{E_{100s}^{(1)}
		=\langle100s|\xi(r)\vec L\cdot\vec S|100s\rangle=0}.
	\]
}






\end{enumerate}
